\documentclass[10pt]{article}

\usepackage[utf8]{inputenc}

\usepackage[T1]{fontenc}

\usepackage{amsmath}

\usepackage{amsfonts}

\usepackage{amssymb}

\usepackage[version=4]{mhchem}

\usepackage{stmaryrd}

\usepackage{graphicx}

\usepackage[export]{adjustbox}

\graphicspath{ {./images/} }



\begin{document}

СВОБОДНАЯ ПЕРЕМЕННАЯ, $\mathrm{c}$ в 0 б 0 д н 0 е в $\mathrm{x}$ 0ж д е н и е п е р е м е н н о й,- вхождение переменной в языковое выражение, являющееся параметром этого выражения. Строгое определение этого понятия может быть дано только для формализованного языка. Для каждого языка дается свое определение С. п., зависящее от правил образования выражений данного языка. Семантич. критерием здесь служит следующее требование: подстановка какого-либо объекта из подразумеваемой интерпретации на место данного вхождения переменной не должна приводить к бессмысленному выражению. Напр., в выражении $\left\{(x, y) \mid x^{2}+y^{2}=z^{2}\right\}$, обозначающем множество точек огружности радиуса $z$, переменная $z$ входит свободно, а переменные $x$ и $y$ нет (см. Связанная переменная). Если $f$ обозначает отображение вида $X \times Y \rightarrow Z$ и переменные $x, y$ пробегают множества $X, Y$ соответственно, то в выражении $f(x, y)$ переменные $x$ и $y$ свободны (так же, как и $f$, если $f$ рассматривать как переменную по функциям). При фиксировании $x$ и варьировании $y$ получается функция вида $Y \rightarrow Z$. Она обозначается через $\lambda y f(x, y)$. В этом выражении $x$ свободно, а $y$ нет. В выражении $(\lambda y f(x, y))(y)$, обозначающем значение функции $\lambda y f(x, y)$ в произвольной точке $y$, последнее вхождение переменной $y$ будет свободным. Два других вхождения не являются свободными. Первое наз. о п е р а т о р н ы м (находящимся под знаком оператора), второе - с в яЗ $\mathbf{2} \mathbf{~ Н ~ Ы ~} \mathbf{~ М . ~}$



Для неформализованных языков, т.е. в реальных математич. текстах, у отдельно взятого выражения не всегда можно определенно выяснить какие переменные у него свободны, а какие связанны. Напр., в выражении $\sum_{i<k} a_{i k}$ в зависимости от контекста переменная $i$ может быть свободной, а $k$ связанной или наоборот, но обе свободными быть не могут. Указание на то, какая переменная считается свободной, дается с помощью дополнительных средств. Напр., если это выражение встречается в контексте вида «пусть $f(k)=$ $=\sum_{i<k} a_{i k}$, то $k$ свободна. Иногда принимают соглашение, что по $k$ суммирование производиться не будет; тогда $k$ - это параметр. Выражение $\left\{a_{i}\right\}$, часто используемое в математике, обозначает иногда одноэлементное множество, п тогда переменная $i$ входит свободно; а иногда оно обозначает множество всех $a_{i}$, когда $i$ пробегает отведенную ему область предметов, и тогда переменная $i$ входит связанно. СВОБОДНАЯ ПОЛУГРУППА н а д а л ф̆ а в и т о м $A$ - полугруппа, әлементами к-рой являются всевозможные конечные последовательности әлементов из $A$ (букв), а операция состоит в приписывании одной последовательности к другой. Элементы С. п. принято называть словами, а операцию часто называют к о н к а т е н а ц и е й. Ради удобства нередко рассматривают также и пустое слово 1 (длина к-рого по определению равна нулю), полагая $w 1=1 w=w$ для любого слова $w$; возникающая таким образом полугруппа с единицей наз. с во бод н м м о но и д о м над $A$. С. П. (свободный моноцд) над $A$ часто обозначают $A+$ (соответственно $\left.\mathrm{A}^{*}\right)$. Для С. П. $A+$ алфоравит $A$ является единственным неприводимым порождающим множеством; он состоит в точности из элементов, неразложимых в произведение. Буквы йз $A$ наз. с в о б о д н ы м и о бр а 3 у щ и м и. С. П. определяется однозначно с точностью до изоморфизма мощностью своего алфавита; эта мощность наз. р а н г о м с в о б о д н о й п о л уг р у п п ы. С. п. ранга 2 имеет подполугруппы, являющиеся С. I. счетного ранга.



С. П. являются свободньции алгео́рами в классе всех полугрупп. Следующие условия для полугруппы $F$ әквивалентны: 1) $F$ есть С. п.; 2) $F$ имеет порождаю- щее множество $\boldsymbol{A}$ такое, что любой элемент из $F$ единственным образом представим в виде произведения әлементов из $A, 3) F$ удовлетворяет закону сокращения, не содержит идемпотентов, каждый әлемент из $F$ имеет конечное число делителей, и для любых $u, v, u^{\prime}, v^{\prime} \in F$ равенство $u v=u^{\prime} v^{\prime}$ влечет, что $u=u^{\prime}$ или один из әлементов $u, u^{\prime}$ есть левый делитель другого.



Всякая подполугруппа $H$ в С. п. имеет единственное неприводимое порождающее множество, состоящее из элементов, неразложимых в $H$ в произведение; однако не всякая подполугруппа С. п. сама свободна. Следующие условия для подполугруппы $H$ в С. п. $F$ әквивалентны: 1) $H$ есть С. п.; 2) для любого $w \in F$ из того, что $w H \cap H \neq \phi$ и $H w \cap H \neq \varnothing$, следует, что $w \in H$; 3) для любого $w \in F$ из того, что $w H \cap H w \cap H \neq \phi$, следует, что $w \in H$. Для произвольных различных слов $u$, $v$ в С. П. $F$ либо $u$ и $v$ являются свободными образующими порожденной ими подполугруппы, либо существует $w \in F$ такое, что $u=w^{k}, v=w^{l}$ для нек-рых натуральных $k, l$; вторая альтернатива выполняется тогда и только тогда, когда $u v=v u$. Всякаяं подполугруппа с тремя образующими в С. п. будет конетно определенной полугруппой, но существуют подполугруппы с четырымя образующими, не являющиеся конечно определенными.



С. П. естественно возникают в автолатов алгебраическай теории (см. также [5], [6]), теории кодирования (см. Кодирование алђавитное, [4] - [6]), теории формальных языков и формальных арамматик (см. [3], [5], [6]). С указанными областями связана проблематика решения уравнений в С. п. (см. [7] - [9]). Существует алгоритм, распознающий разрешимость произвольных уравнений в С. П.



Лит.: [1] К ли ф ф о р д А., П р е с т о н Г., Алгебраическая теория полугрупп, пер. с англ., т. $1-2$, М., 1972; [2] Л яп и н Е. С., Полугруппы, М., $1960 ;$ [3] Гро с с М., Л а нт е н А., Теория формальных грамматик, пер. с франц., М., $1971 ;$ [4] М а p к о в A. А., Введение в теорию кодирования, M., 1982; [5] E i l e n b e r g S., Automata, languages and machines, v. A-B, N. Y.-L., $1974-76$; [6] L a 1 l e m e n t G. chines, V. A-B, N. Y.L., $1974-76 ;[6]$ L a 11 e m e n t G.j,

Semigroups and combinatorial applications, N. Y. [a. o.], 1979; [7] L e n t i n A., Equations dans les monoïdes libres, $1979 ;$ [7] L e n t i n A., Equations dans les monoïdes libres,

Р., 1972; [8] Х м е л е в с к и й Ю. И., Уравнения в свободной полугруппе, М., 1971 (Тр. Матем. ин-та АН СССР, т. 107) [9] М а к а н и н Г. С., «Матем. сб.», 1977, т. 103, № 2, с. 147 СВОБОДНАЯ РЕЗОЛБВЕНТА - частный случай проективной резольвенты. Всякий модуль $M$ над ассоциативным кольцом $R$ является фактормодулем $F_{0} / N_{0}$ свободного $R$-модуля $F_{0}$ по нек-рому его подмодулю $N_{0}$. Для подмодуля $N_{0}$ существует аналогичное представление $F_{1} / N_{1}$ и т. Д. В результате получается точная последовательность свободных модулей



\[

F_{0} \leftarrow F_{1} \leftarrow F_{2} \leftarrow \ldots \leftarrow F_{n} \leftarrow

\]



называемая С. р. модуля $M$. Канонич. гомоморфиззм



\includegraphics[max width=\textwidth, center]{2023_07_08_32e26af960ac4ff6da9bg-1}

$\begin{array}{ll}\text { М о м. } & \text { В. Е. Говоров }\end{array}$



СВОБОДНАЯ РЕIIЕТКА - свободная алаебра многообразия всех решеток. Решены [1] проблемы тождества слов и канонич. представления слова в С. р.



Jum.: [1] W h i t m a n P. M., "Annals Math.", 1941, v. 42,

$325-30 ; 1942$, v. 43 , p. $104-15$.



СВОБОДНО СТАНОВЯШАЯСЯ ПОСЛЕДОВАТЕЛЬНОСТБ - см. Интуиционизм.



СВОБОДНОЕ ГАРМОНИЧЕСКОЕ КОЛЕБАНИЕ синусоидальное колебание. Если механическая или физич. величина $x(t)$, где $t$ - время, меняется по закону



\[

x(t)=A \cos (\omega t+\varphi)

\]



то говорят, что $x(t)$ совершает С. г. к. Здесь $A, \omega, \varphi-$ действительные постоянные, $A>0, \omega>0$. Величины $A, \omega, \varphi$ наз. соответственно а м п л и т у Д о й, ч а сто той, ф а 3 ой С. г.к. П е ри и д С. г. к. равен





\end{document}